\section{Introduction}
\label{sec:intro}

Halal compliance is, at the regulatory frontier, a qualitative discipline.
Standards such as MS~1500:2019 \cite{ms1500}, OIC/SMIIC~1 \cite{smiic1}, and
GSO~2055-1 \cite{gso2055} encode juristic principles into operational
checklists, but the binding judgments --- ``reasonable precaution'' against
cross-contamination, ``trustworthy'' certification chains, ``negligible''
trace contaminants --- remain expert-mediated and locally interpreted. The
result is a globally fragmented compliance landscape where mutual
recognition between certifying bodies is the exception rather than the
rule, where supply-chain shocks expose the limits of paper-based
attestations, and where billions of consumers cannot verify, even in
principle, the provenance claims attached to the products they buy.

This paper introduces \textbf{Global Halal Compliance (GHC)}, a
cross-domain protocol that replaces qualitative judgments with verified,
quantitative instruments without displacing the underlying juristic
authority. GHC's central thesis is that the compositional structure of a
modern supply chain --- batches, transformations, mixings, separations,
audits --- is exactly the structure studied by category theory
\cite{maclane,fongspivak}, and that the propagation of compliance status
through that structure is naturally modeled by the tools of information
theory \cite{covertom}, open-system quantum mechanics \cite{nielsenchuang},
spectral graph theory \cite{chung}, and zero-knowledge cryptography
\cite{gmr1989,groth16}.

\paragraph{Contributions.}
The remainder of this paper develops six contributions, each
machine-checked in Lean~4 \cite{lean4}, implemented in the GHC reference
codebase, and standardized in the GHC Protocol specification:
\begin{enumerate}
  \item The \textbf{Halal-Closure Functor} $\HCF \colon \SC \to \Lat$
        (\Cref{sec:hcf}), establishing compositional semantics for
        \halal/\mashbuh/\haram\ status.
  \item The \textbf{Contamination Channel Capacity} $\Ch$
        (\Cref{sec:capacity}), a closed-form upper bound on haram-signal
        survival.
  \item The \textbf{Najis Density Matrix} $\rhoN$
        (\Cref{sec:najis}), giving a single scalar provenance-purity
        score.
  \item \textbf{Spectral provenance invariants} (\Cref{sec:spectral})
        for tamper-evident DAG fingerprinting.
  \item The \textbf{zk-Halal} attestation scheme (\Cref{sec:zkhalal})
        reconciling transparency with supplier privacy.
  \item The \textbf{Compliance Convergence Theorem}
        (\Cref{sec:convergence}) quantifying audit sufficiency.
\end{enumerate}

\paragraph{Scope and non-goals.}
GHC formalizes the \emph{propagation} of compliance status; it does not
adjudicate the underlying juristic questions. Edge cases (e.g.\ the
status of a particular emulsifier, or the validity of a specific
slaughter method) remain the prerogative of recognized Shariah
authorities. GHC's compliance lattice is parametric in the authority,
and the protocol explicitly supports federated multi-authority operation
with first-class encoding of dissent.
